\chapter{Απαιτήσεις και σχεδίαση συστήματος}
\label{ch:design}

\guide{Έκταση 12--15 σελίδες. Καλό σημείο εκκίνησης για τη γραφή: το σύστημα το ξέρεις ήδη. Περιγράφεις \emph{τι} και \emph{γιατί}, όχι \emph{πώς} (αυτό είναι το \cref{ch:implementation}).}

\section{Σενάρια χρήσης}
\label{sec:usecases}

\guide{Τρία σενάρια, το καθένα σε μία παράγραφο με τα βήματα του χρήστη:
(α) \textbf{Επιθεώρηση πεδίου}: φόρτωση έργου από Terraweb, ευθυγράμμιση σε 2--4 σημεία ελέγχου, περιήγηση στους αισθητήρες, συναγερμοί, διαδρομή επιθεώρησης, αναφορά.
(β) \textbf{Προετοιμασία στο γραφείο}: το μοντέλο ως μακέτα πάνω στο τραπέζι, χρονομηχανή, χάρτης θερμότητας.
(γ) \textbf{Εκπαίδευση}: το εικονικό έργο εξάσκησης (\code{PracticeProject}) και ο οδηγός πρώτης χρήσης (\code{Onboarding}).}

\section{Απαιτήσεις}
\label{sec:requirements}

\subsection{Λειτουργικές απαιτήσεις}

\begin{table}[htbp]
  \centering\small
  \caption{Λειτουργικές απαιτήσεις.}
  \label{tab:functional-req}
  \begin{tabularx}{\textwidth}{@{}lX@{}}
    \toprule
    Κωδ. & Απαίτηση \\
    \midrule
    ΛΑ1 & Λήψη έργου (σημεία αισθητήρων με συντεταγμένες κανάβου) από την πλατφόρμα παρακολούθησης ή από αρχείο CSV. \\
    ΛΑ2 & Ευθυγράμμιση του έργου με τον πραγματικό χώρο με άγγιγμα γνωστών σημείων και επιβεβαίωση από τον χρήστη. \\
    ΛΑ3 & Διατήρηση της ευθυγράμμισης μεταξύ συνεδριών. \\
    ΛΑ4 & Απεικόνιση κάθε αισθητήρα στη θέση του με τύπο, όνομα, κατάσταση και τελευταία τιμή. \\
    ΛΑ5 & Απεικόνιση ιστορικού μετρήσεων και κατάστασης σε παρελθούσα ημερομηνία. \\
    ΛΑ6 & Ειδοποίηση για αλλαγές κατάστασης συναγερμού σε (σχεδόν) πραγματικό χρόνο. \\
    ΛΑ7 & Υπέρθεση μοντέλου 3Δ (OBJ, DXF, IFC) στο ίδιο σύστημα συντεταγμένων. \\
    ΛΑ8 & Καθοδηγούμενη διαδρομή επιθεώρησης με αναφορά. \\
    ΛΑ9 & Πρόσβαση τρίτων εφαρμογών μέσω τεκμηριωμένου API. \\
    \bottomrule
  \end{tabularx}
\end{table}

\subsection{Μη λειτουργικές απαιτήσεις}

\guide{ΜΛΑ: σταθερός ρυθμός καρέ (72\,Hz), λειτουργία χωρίς σύνδεση (\code{OfflineCache}), ασφάλεια σε τοπικό δίκτυο (κλειδί API, όρια μεγέθους, rate limiting), χρήση με λεπτά γάντια εργασίας (μεγάλοι στόχοι, κουμπιά αντί για χειρονομίες), δίγλωσση διεπαφή (ελληνικά/αγγλικά), ανεξαρτησία από συγκεκριμένο CRS, άνεση (αποφυγή ναυτίας στην κίνηση με stick). Για καθεμία: γιατί προκύπτει από το πεδίο.}

\section{Αρχιτεκτονική}
\label{sec:architecture}

Το σύστημα αποτελείται από τρία επίπεδα (\cref{fig:architecture}): την πλατφόρμα παρακολούθησης στο νέφος, προαιρετικά εργαλεία σε υπολογιστή στο τοπικό δίκτυο, και την εφαρμογή στη συσκευή, η οποία λειτουργεί και αυτόνομα.

\begin{figure}[htbp]
  \centering
  \resizebox{\textwidth}{!}{% Διάγραμμα αρχιτεκτονικής (TikZ). Χρήση: % Διάγραμμα αρχιτεκτονικής (TikZ). Χρήση: % Διάγραμμα αρχιτεκτονικής (TikZ). Χρήση: \input{../figures/architecture} μέσα σε figure.
% Τα ονόματα κλάσεων/αρχείων είναι τα πραγματικά του QuestMarkerApp.
\begin{tikzpicture}[
  font=\sffamily\scriptsize,
  node distance=4mm and 6mm,
  comp/.style={draw, rounded corners=2pt, fill=white, align=center, minimum width=27mm, minimum height=7mm, inner sep=2pt},
  tier/.style={draw=#1!60!black, fill=#1!8, rounded corners=4pt, inner sep=3mm},
  tiertitle/.style={font=\sffamily\small\bfseries, anchor=south west},
  flow/.style={-{Stealth[length=2mm]}, thick},
  lbl/.style={font=\sffamily\tiny, fill=white, inner sep=1pt, align=center},
]

% ---------------- Cloud -------------------------------------------------
\node[comp] (tw) {Terraweb\\(REST, JWT)};
\node[comp, right=of tw] (wit) {Wit.ai\\(φωνή, προαιρετικό)};

% ---------------- PC ----------------------------------------------------
\node[comp, below=14mm of tw, xshift=-16mm] (bridge) {\texttt{terraweb\_bridge.py}\\\texttt{send\_project.py}};
\node[comp, right=of bridge] (models) {\texttt{send\_model.py}\\\texttt{model\_prep.py}};
\node[comp, right=of models] (cserver) {\texttt{coordinate\_server.py}\\(long-poll \texttt{/api/events})};

% ---------------- Quest: network ------------------------------------------
\node[comp, below=16mm of bridge] (api) {\texttt{LocalAPIServer}\\HTTP :8080};
\node[comp, right=of api] (twc) {\texttt{TerrawebClient}\\\texttt{LiveAlerts}};
\node[comp, right=of twc] (push) {\texttt{WebSocketPushClient}\\\texttt{NetworkDiscovery}};
\node[comp, below=3mm of api, minimum width=27mm] (sec) {\texttt{ApiSecurity}\\κλειδί, rate limit};

% ---------------- Quest: core ---------------------------------------------
\node[comp, below=8mm of sec] (align) {\texttt{ProjectAlignment}\\\texttt{CoordinateTransform}};
\node[comp, right=of align] (anch) {\texttt{SpatialAnchorManager}\\\texttt{AnchorFieldTest}};
\node[comp, right=of anch] (sess) {\texttt{SessionManager}\\\texttt{OfflineCache}};

% ---------------- Quest: visualization / UI -------------------------------
\node[comp, below=of align] (mark) {\texttt{MarkerManager}\\\texttt{HistoryRibbons}};
\node[comp, right=of mark] (model) {\texttt{ModelOverlay}\\\texttt{ModelHeatmap}};
\node[comp, right=of model] (ui) {\texttt{MainMenu}, \texttt{SensorCard}\\\texttt{VoiceCommands}, \texttt{InspectionRoute}};

% ---------------- Tiers ---------------------------------------------------
\begin{scope}[on background layer]
  \node[tier=gray, fit=(tw)(wit)] (cloudT) {};
  \node[tier=green, fit=(bridge)(models)(cserver)] (pcT) {};
  \node[tier=blue, fit=(api)(push)(sec)(align)(sess)(mark)(ui)] (questT) {};
\end{scope}
\node[tiertitle] at (cloudT.north west) {Νέφος};
\node[tiertitle] at (pcT.north west) {PC στο LAN (Python)};
\node[tiertitle] at (questT.north west) {Meta Quest 3 (Unity, C\#)};

% ---------------- Flows ---------------------------------------------------
\draw[flow] (tw) -- node[lbl]{αισθητήρες} (bridge);
\draw[flow] (bridge) -- node[lbl]{POST /api/project} (api);
\draw[flow] (models) -- node[lbl, pos=0.35]{POST /api/model\\(zip)} (api);
\draw[flow] (cserver) -- node[lbl]{push} (push);
\draw[flow, <->] (tw.east) -| ++(4mm,-20mm) -| node[lbl, pos=0.75]{HTTPS: μετρήσεις,\\ιστορικό, συναγερμοί} (twc.north);
\draw[flow, <->, dashed] (wit.south) |- ($(ui.north east)+(3mm,8mm)$) -- ($(ui.north east)+(3mm,0)$);
\draw[flow] (api) -- (sec);
\draw[flow] (sec) -- (align);
\draw[flow] (align) -- (mark);
\draw[flow] (align) -- (anch);
\draw[flow] (align.south east) -- (model);
\draw[flow] (twc) -- ++(0,-8mm) -| (sess);
\end{tikzpicture}
 μέσα σε figure.
% Τα ονόματα κλάσεων/αρχείων είναι τα πραγματικά του QuestMarkerApp.
\begin{tikzpicture}[
  font=\sffamily\scriptsize,
  node distance=4mm and 6mm,
  comp/.style={draw, rounded corners=2pt, fill=white, align=center, minimum width=27mm, minimum height=7mm, inner sep=2pt},
  tier/.style={draw=#1!60!black, fill=#1!8, rounded corners=4pt, inner sep=3mm},
  tiertitle/.style={font=\sffamily\small\bfseries, anchor=south west},
  flow/.style={-{Stealth[length=2mm]}, thick},
  lbl/.style={font=\sffamily\tiny, fill=white, inner sep=1pt, align=center},
]

% ---------------- Cloud -------------------------------------------------
\node[comp] (tw) {Terraweb\\(REST, JWT)};
\node[comp, right=of tw] (wit) {Wit.ai\\(φωνή, προαιρετικό)};

% ---------------- PC ----------------------------------------------------
\node[comp, below=14mm of tw, xshift=-16mm] (bridge) {\texttt{terraweb\_bridge.py}\\\texttt{send\_project.py}};
\node[comp, right=of bridge] (models) {\texttt{send\_model.py}\\\texttt{model\_prep.py}};
\node[comp, right=of models] (cserver) {\texttt{coordinate\_server.py}\\(long-poll \texttt{/api/events})};

% ---------------- Quest: network ------------------------------------------
\node[comp, below=16mm of bridge] (api) {\texttt{LocalAPIServer}\\HTTP :8080};
\node[comp, right=of api] (twc) {\texttt{TerrawebClient}\\\texttt{LiveAlerts}};
\node[comp, right=of twc] (push) {\texttt{WebSocketPushClient}\\\texttt{NetworkDiscovery}};
\node[comp, below=3mm of api, minimum width=27mm] (sec) {\texttt{ApiSecurity}\\κλειδί, rate limit};

% ---------------- Quest: core ---------------------------------------------
\node[comp, below=8mm of sec] (align) {\texttt{ProjectAlignment}\\\texttt{CoordinateTransform}};
\node[comp, right=of align] (anch) {\texttt{SpatialAnchorManager}\\\texttt{AnchorFieldTest}};
\node[comp, right=of anch] (sess) {\texttt{SessionManager}\\\texttt{OfflineCache}};

% ---------------- Quest: visualization / UI -------------------------------
\node[comp, below=of align] (mark) {\texttt{MarkerManager}\\\texttt{HistoryRibbons}};
\node[comp, right=of mark] (model) {\texttt{ModelOverlay}\\\texttt{ModelHeatmap}};
\node[comp, right=of model] (ui) {\texttt{MainMenu}, \texttt{SensorCard}\\\texttt{VoiceCommands}, \texttt{InspectionRoute}};

% ---------------- Tiers ---------------------------------------------------
\begin{scope}[on background layer]
  \node[tier=gray, fit=(tw)(wit)] (cloudT) {};
  \node[tier=green, fit=(bridge)(models)(cserver)] (pcT) {};
  \node[tier=blue, fit=(api)(push)(sec)(align)(sess)(mark)(ui)] (questT) {};
\end{scope}
\node[tiertitle] at (cloudT.north west) {Νέφος};
\node[tiertitle] at (pcT.north west) {PC στο LAN (Python)};
\node[tiertitle] at (questT.north west) {Meta Quest 3 (Unity, C\#)};

% ---------------- Flows ---------------------------------------------------
\draw[flow] (tw) -- node[lbl]{αισθητήρες} (bridge);
\draw[flow] (bridge) -- node[lbl]{POST /api/project} (api);
\draw[flow] (models) -- node[lbl, pos=0.35]{POST /api/model\\(zip)} (api);
\draw[flow] (cserver) -- node[lbl]{push} (push);
\draw[flow, <->] (tw.east) -| ++(4mm,-20mm) -| node[lbl, pos=0.75]{HTTPS: μετρήσεις,\\ιστορικό, συναγερμοί} (twc.north);
\draw[flow, <->, dashed] (wit.south) |- ($(ui.north east)+(3mm,8mm)$) -- ($(ui.north east)+(3mm,0)$);
\draw[flow] (api) -- (sec);
\draw[flow] (sec) -- (align);
\draw[flow] (align) -- (mark);
\draw[flow] (align) -- (anch);
\draw[flow] (align.south east) -- (model);
\draw[flow] (twc) -- ++(0,-8mm) -| (sess);
\end{tikzpicture}
 μέσα σε figure.
% Τα ονόματα κλάσεων/αρχείων είναι τα πραγματικά του QuestMarkerApp.
\begin{tikzpicture}[
  font=\sffamily\scriptsize,
  node distance=4mm and 6mm,
  comp/.style={draw, rounded corners=2pt, fill=white, align=center, minimum width=27mm, minimum height=7mm, inner sep=2pt},
  tier/.style={draw=#1!60!black, fill=#1!8, rounded corners=4pt, inner sep=3mm},
  tiertitle/.style={font=\sffamily\small\bfseries, anchor=south west},
  flow/.style={-{Stealth[length=2mm]}, thick},
  lbl/.style={font=\sffamily\tiny, fill=white, inner sep=1pt, align=center},
]

% ---------------- Cloud -------------------------------------------------
\node[comp] (tw) {Terraweb\\(REST, JWT)};
\node[comp, right=of tw] (wit) {Wit.ai\\(φωνή, προαιρετικό)};

% ---------------- PC ----------------------------------------------------
\node[comp, below=14mm of tw, xshift=-16mm] (bridge) {\texttt{terraweb\_bridge.py}\\\texttt{send\_project.py}};
\node[comp, right=of bridge] (models) {\texttt{send\_model.py}\\\texttt{model\_prep.py}};
\node[comp, right=of models] (cserver) {\texttt{coordinate\_server.py}\\(long-poll \texttt{/api/events})};

% ---------------- Quest: network ------------------------------------------
\node[comp, below=16mm of bridge] (api) {\texttt{LocalAPIServer}\\HTTP :8080};
\node[comp, right=of api] (twc) {\texttt{TerrawebClient}\\\texttt{LiveAlerts}};
\node[comp, right=of twc] (push) {\texttt{WebSocketPushClient}\\\texttt{NetworkDiscovery}};
\node[comp, below=3mm of api, minimum width=27mm] (sec) {\texttt{ApiSecurity}\\κλειδί, rate limit};

% ---------------- Quest: core ---------------------------------------------
\node[comp, below=8mm of sec] (align) {\texttt{ProjectAlignment}\\\texttt{CoordinateTransform}};
\node[comp, right=of align] (anch) {\texttt{SpatialAnchorManager}\\\texttt{AnchorFieldTest}};
\node[comp, right=of anch] (sess) {\texttt{SessionManager}\\\texttt{OfflineCache}};

% ---------------- Quest: visualization / UI -------------------------------
\node[comp, below=of align] (mark) {\texttt{MarkerManager}\\\texttt{HistoryRibbons}};
\node[comp, right=of mark] (model) {\texttt{ModelOverlay}\\\texttt{ModelHeatmap}};
\node[comp, right=of model] (ui) {\texttt{MainMenu}, \texttt{SensorCard}\\\texttt{VoiceCommands}, \texttt{InspectionRoute}};

% ---------------- Tiers ---------------------------------------------------
\begin{scope}[on background layer]
  \node[tier=gray, fit=(tw)(wit)] (cloudT) {};
  \node[tier=green, fit=(bridge)(models)(cserver)] (pcT) {};
  \node[tier=blue, fit=(api)(push)(sec)(align)(sess)(mark)(ui)] (questT) {};
\end{scope}
\node[tiertitle] at (cloudT.north west) {Νέφος};
\node[tiertitle] at (pcT.north west) {PC στο LAN (Python)};
\node[tiertitle] at (questT.north west) {Meta Quest 3 (Unity, C\#)};

% ---------------- Flows ---------------------------------------------------
\draw[flow] (tw) -- node[lbl]{αισθητήρες} (bridge);
\draw[flow] (bridge) -- node[lbl]{POST /api/project} (api);
\draw[flow] (models) -- node[lbl, pos=0.35]{POST /api/model\\(zip)} (api);
\draw[flow] (cserver) -- node[lbl]{push} (push);
\draw[flow, <->] (tw.east) -| ++(4mm,-20mm) -| node[lbl, pos=0.75]{HTTPS: μετρήσεις,\\ιστορικό, συναγερμοί} (twc.north);
\draw[flow, <->, dashed] (wit.south) |- ($(ui.north east)+(3mm,8mm)$) -- ($(ui.north east)+(3mm,0)$);
\draw[flow] (api) -- (sec);
\draw[flow] (sec) -- (align);
\draw[flow] (align) -- (mark);
\draw[flow] (align) -- (anch);
\draw[flow] (align.south east) -- (model);
\draw[flow] (twc) -- ++(0,-8mm) -| (sess);
\end{tikzpicture}
}
  \caption{Αρχιτεκτονική του συστήματος.}
  \label{fig:architecture}
\end{figure}

\guide{Περιέγραψε κάθε επίπεδο και κάθε ροή του σχήματος. Τόνισε την επιλογή: ο \emph{εξυπηρετητής} τρέχει στα γυαλιά (αντί η εφαρμογή να είναι μόνο πελάτης), ώστε οποιαδήποτε συσκευή στο LAN να στέλνει δεδομένα χωρίς ενδιάμεσο. Και ότι η εφαρμογή μιλά και απευθείας με το Terraweb, χωρίς PC.}

\section{Επικοινωνία}
\label{sec:communication}

\subsection{Τοπικό REST API}

\Cref{tab:api-summary} συνοψίζει τα endpoints του ενσωματωμένου εξυπηρετητή (θύρα 8080). Η πλήρης αναφορά βρίσκεται στο \cref{app:api}.

\begin{table}[htbp]
  \centering\small
  \caption{Ομάδες endpoints του τοπικού API.}
  \label{tab:api-summary}
  \begin{tabularx}{\textwidth}{@{}lX@{}}
    \toprule
    Ομάδα & Endpoints \\
    \midrule
    Σημεία & \code{/api/coordinates}, \code{/api/clear}, \code{/api/undo}, \code{/api/layer} \\
    Έργο / ευθυγράμμιση & \code{/api/project}, \code{/api/project/reset}, \code{/api/calibrate/*}, \code{/api/gps-origin}, \code{/api/anchors} \\
    Μοντέλα & \code{/api/model}, \code{/api/models} \\
    Συνεδρίες & \code{/api/session}, \code{/api/session/new} \\
    Πεδίο & \code{/api/fieldnotes}, \code{/api/reports}, \code{/api/fieldtest} \\
    Διάγνωση & \code{/api/health}, \code{/api/status}, \code{/api/logs} \\
    \bottomrule
  \end{tabularx}
\end{table}

\subsection{Ανακάλυψη και push}
\guide{\code{NetworkDiscovery}: εκπομπή UDP στη θύρα 48080 με JSON \{service, port, name\}, ώστε τα εργαλεία PC να βρίσκουν τα γυαλιά χωρίς IP. Push μέσω long-poll στο \code{/api/events} του \code{coordinate_server.py} (\code{WebSocketPushClient}). Γιατί long-poll και όχι WebSocket/MQTT (απλότητα, διέλευση από firewalls, καμία εξάρτηση στη συσκευή).}

\subsection{Σύνδεση με την πλατφόρμα Terraweb}
\guide{Αυθεντικοποίηση JWT (access 15 λεπτά, refresh 7 ημέρες), επιλογή έργου, λήψη αισθητήρων, ιστορικού και ενεργών συναγερμών (polling ανά 20\,s). Μόνο HTTPS για διαπιστευτήρια· pinning πιστοποιητικού για αυτο-υπογεγραμμένους εξυπηρετητές.}

\section{Ασφάλεια}
\label{sec:security}
\guide{Μοντέλο απειλών για LAN: κλειδί API (\code{X-Api-Key} ή \code{Bearer}), όριο ρυθμού αιτημάτων, όρια μεγέθους σώματος, προστασία από zip bomb στα πακέτα μοντέλων, καμία έκθεση εκτός LAN. Τι \emph{δεν} καλύπτεται (cleartext HTTP εντός LAN) και γιατί είναι αποδεκτό εδώ.}

\section{Μοντέλο δεδομένων}
\label{sec:datamodel}
\guide{\code{CoordinateData} (σημείο: id, x/y/z, χρώμα, ετικέτα, στρώμα), \code{ProjectData} (σημεία έργου με id, όνομα, τύπο, E/N/H), συνεδρία, εγγραφές ιστορικού. Σχήμα κλάσεων UML με τα κύρια πεδία.}
